\section{Modellering}

\subsection{Fra værhypotese til frozen candidate}

Utviklingen fulgte en kronologisk ML lifecycle med værhypotesen som utgangspunkt.
Målet var å undersøke om weather features kunne gi lavere feil enn modeller basert
på historisk trafikk. Først ble forrige-uke-prediksjonen etablert som kontroll. For
en forecast med target $t+1$ bruker baselinen observasjonen ved
$(t+1)-168=t-167$, altså samme target-time én uke tidligere. En lineær regresjon og
en \texttt{HistGradientBoostingRegressor} (HGB) ble deretter trent med kalender og
traffic lag features ved $t-1$, $t-2$, $t-24$ og $t-168$. HGB er en
histogram-based gradient boosting model for tabulære data. På validation i 2024
fikk lineær regresjon MAE 322,45, HGB 103,28 og baselinen 200,78. HGB ble derfor
valgt og oppnådde MAE 106,22 på 2025-testen.

Sammenligningen brukte 8 231 felles rader i 2024 og 8 760 i 2025. Modellen ble først
tilpasset på 2022--2023 for validation. Etter valg ble HGB trent på training- og
validation-data samlet før den ene målingen på 2025-testen. Denne rekkefølgen skiller
mellom model selection og å anslå generalization error. Dersom 2025 hadde blitt brukt til
å endre features etter hver måling, ville perioden gradvis blitt en ny
validation period og sluttresultatet for optimistisk.

Videre analyse viste at store feil ofte sammenfalt med helligdager og høytider.
Deterministiske norske kalenderindikatorer ble derfor vurdert med tre expanding
folds i 2022--2023. Gjennomsnittlig CV-MAE ble 119,69 med indikatorene og 120,03
uten. Deretter ble den opprinnelige værhypotesen testet direkte. En raw weather model
tapte mot traffic-only model på felles 2024-rader, med MAE 106,73 mot 103,41. En
forhåndsdefinert weather feature-pakke med blant annet akkumulert nedbør og
sesongavvik for temperatur og vind tapte senere alle tre CV-foldene: 124,43 mot
119,69. Vær ga altså ingen dokumentert forbedring i noen av de to forsøkene og ble
forkastet før sluttmodellen ble bestemt.

Den første frozen calendar HGB-en oppnådde MAE 102,35 på 2025, men krevde de ferske
lag features $t-1$ og $t-2$. Traffic Data API kan publisere timeaggregater med forsinkelse,
så modellen var lite egnet til en reell neste-time-forespørsel. Prosjektet definerte
derfor tre små, latency-robust lagpakker på forhånd: $(24,168)$, $(24,48,168)$ og
$(24,168,336)$. Alle brukte de samme calendar feature-typene og HGB med fast
\texttt{random\_state=42}. Gjennomsnittlig CV-MAE i 2022--2023 ble henholdsvis
134,88, 133,77 og 135,80. Pakken $(24,48,168)$ ble fryst som
\texttt{traffic-hgb-latency-robust-v1}.

Den latency-robust modellen er ikke kandidaten med lavest tidligere 2025-MAE; den
løser et annet programvarekrav. Ved å fjerne $t-1$ og $t-2$ kan appen lage en
prognose selv om de nyeste aggregatene ikke er publisert. Kostnaden er at modellen
mister informasjon om den umiddelbare trafikkutviklingen. Kandidatlisten ble holdt
liten for å begrense tilpasning til CV-foldene, og ingen hyperparameters ble valgt
fra 2026-resultatet.

\subsection{Kronologisk evaluering}

Random splitting ville blandet fortid og fremtid og gitt en misvisende vurdering av
tidsserien. Prosjektet brukte derfor time-series cross-validation der hver validation
period ligger etter training period. Model selection brukte expanding training
windows med validation i 2022 H2, 2023 H1 og 2023 H2. Perioden 2025 var allerede
brukt som test set for en tidligere modell og ble derfor låst for videre valg. Etter
at den latency-robust kandidaten var frozen, ble den trent på 2022--2025 og evaluert
én gang på det nye test set-et 2026 H1. Sluttresultatet ble ikke brukt til ny tuning.

Tabell~\ref{tab:final-results} sammenligner modell og baseline på nøyaktig de samme
4 343 radene. Modellen reduserer MAE med 102,98 kjøretøy per time, tilsvarende 48,6
prosent relativt til baselinen. RMSE er også vesentlig lavere, men ligger over MAE
fordi enkelte timer har store feil.

\begin{table}[htbp]
  \centering
  \caption{Sluttevaluering på den tidligere ubrukte perioden 2026 H1.}
  \label{tab:final-results}
  \begin{tabularx}{\textwidth}{Xrrr}
    \toprule
    Modell & MAE & RMSE & Rader \\
    \midrule
    Latensrobust HGB & 108,81 & 161,85 & 4 343 \\
    Forrige uke & 211,79 & 454,47 & 4 343 \\
    \bottomrule
  \end{tabularx}
\end{table}

\subsection{Feature importance}

For å undersøke hva den valgte modellen faktisk bygger prediction på, ble
permutation feature importance beregnet med MAE som scoring. Modellen ble trent på
nytt i hver av de tre kronologiske CV-foldene, og hver feature ble permutert 30
ganger på foldens validation-data. Figur~\ref{fig:feature-importance} viser
gjennomsnittlig økning i MAE; en høy positiv verdi betyr at modellen mister mye
predictive performance når sammenhengen til target brytes. Analysen ble utført etter
at kandidaten var frozen og ble ikke brukt til å endre modell eller features.

\begin{figure}[htbp]
  \centering
  \begin{tikzpicture}
    \begin{axis}[
      xbar,
      width=0.82\textwidth,
      height=8.4cm,
      xlabel={Økning i MAE (kjøretøy/time)},
      table/col sep=comma,
      ytick=data,
      yticklabels from table={figures/feature_importance.csv}{label},
      y dir=reverse,
      xmajorgrids=true,
      grid style={black!12},
      axis line style={black!55},
      tick style={black!55},
      bar width=5pt,
      enlarge y limits=0.035,
      error bars/x dir=both,
      error bars/x explicit
    ]
      \addplot[fill=reportblue!75, draw=reportblue]
        table[
          x=mean_delta_mae,
          x error=std_delta_mae,
          y expr=\coordindex,
          col sep=comma
        ]{figures/feature_importance.csv};
    \end{axis}
  \end{tikzpicture}
  \caption{Cross-validated permutation feature importance for den frozen
  latency-robust modellen. Error bars viser standardavvik over 90 permutasjoner
  per feature, fordelt på tre validation-folds.}
  \label{fig:feature-importance}
\end{figure}

Størst importance fikk \texttt{traffic\_168h\_ago} med en gjennomsnittlig
MAE-økning på 394,64, etterfulgt av \texttt{hour\_cos} med 366,88,
\texttt{traffic\_24h\_ago} med 223,15 og \texttt{day\_of\_week} med 182,21.
Resultatet viser at modellen særlig støtter seg på ukerytme, døgnsyklus og trafikken
ett døgn tidligere.

Resultatet må tolkes som model reliance, ikke som en kausal effekt. Flere inputs
deler informasjon: \texttt{hour}, \texttt{hour\_sin} og \texttt{hour\_cos} beskriver
samme døgnsyklus, mens traffic lag features er korrelerte. Permutation importance
kan derfor fordele eller undervurdere importance mellom slike features. Analysen
gjelder dessuten bare features i sluttmodellen; konklusjonen om at vær ikke hjalp
kommer fortsatt fra de separate weather-model-sammenligningene.

\subsection{Feilanalyse}

Figur~\ref{fig:hourly-error} viser MAE etter target-time i 2026 H1. Feilen er lavest
kl. 04 med 32,24 og høyest kl. 17 med 152,58 kjøretøy per time. Perioden fra kl. 07
til 18 er gjennomgående vanskeligere enn natten, noe som er rimelig når volumet og
endringstakten er større. Kl. 21 har også høy MAE, 151,36, og viser at utfordringen
ikke bare kan reduseres til klassisk rushtrafikk.

\begin{figure}[htbp]
  \centering
  \begin{tikzpicture}
    \begin{axis}[
      width=0.95\textwidth,
      height=6.2cm,
      xlabel={Target-time i Europe/Oslo},
      ylabel={MAE kjøretøy/time},
      xmin=-0.5, xmax=23.5,
      ymin=0, ymax=175,
      xtick={0,3,6,9,12,15,18,21,23},
      ymajorgrids=true,
      grid style={black!12},
      axis line style={black!55},
      tick style={black!55}
    ]
      \addplot[reportblue, semithick, mark=*, mark size=1.8pt]
        table[x=hour,y=mae,col sep=comma]{figures/error_by_hour.csv};
    \end{axis}
  \end{tikzpicture}
  \caption{Sluttmodellens MAE per target-time i 2026 H1. Figuren genereres fra
  prosjektets dokumenterte sluttrapport.}
  \label{fig:hourly-error}
\end{figure}

Etter ukedag varierer MAE fra 97,05 på tirsdag til 121,89 på lørdag. De øvrige
dagene ligger mellom 100,14 og 112,98. Analysen viser når feilen er størst, men
forklarer ikke årsaken alene. Aktuelle forbedringer må testes i en ny kronologisk
validation-runde, for eksempel hendelsesdata eller bedre mål på
publiseringsforsinkelse, før en ny og senere testperiode åpnes.

RMSE 161,85 mot MAE 108,81 viser at fordelingen inneholder noen større avvik. For
en bruker betyr dette at gjennomsnittsfeilen ikke bør leses som en garanti for en
enkelt time. Appen viser derfor et tallfestet anslag og historikk, men ingen
sikkerhetskategori som kunne gitt et falskt inntrykk av kalibrert usikkerhet.
